\documentclass[11pt]{article}

\usepackage[margin=1in]{geometry}
\usepackage{times}
\usepackage{microtype}
\usepackage{amsmath,amssymb}
\usepackage{booktabs}
\usepackage{graphicx}
\usepackage{xcolor}
\usepackage{natbib}
\usepackage{algorithm}
\usepackage{algpseudocode}
\usepackage{caption}
\usepackage[colorlinks=true,allcolors=black!70]{hyperref}

\bibliographystyle{plainnat}

% Placeholders are deliberately loud. Nothing in this paper may be mistaken for a
% result until a result exists.
\definecolor{gapbg}{RGB}{255,244,230}
\definecolor{gaprule}{RGB}{200,120,0}
\newcommand{\gap}[1]{%
  \par\vspace{0.6em}\noindent
  \colorbox{gapbg}{%
    \parbox{\dimexpr\linewidth-2\fboxsep}{%
      \color{gaprule}\rule{\linewidth}{0.8pt}\\[0.2em]
      \textcolor{gaprule}{\textbf{[PENDING EXPERIMENT]}}\ \color{black}#1\\[0.3em]
      \color{gaprule}\rule{\linewidth}{0.8pt}}}%
  \par\vspace{0.6em}}

\title{\bf When Does Deep Forecasting Stop Paying?\\
A Benchmark for Demand Forecasting Under Data Scarcity}

\author{Aryan Shirodkar\\
\normalsize \textit{Independent Researcher}\\
\normalsize \texttt{github.com/aryanshirodkar15/scarce-forecasting}}

\date{\today}

\begin{document}
\maketitle

\begin{abstract}
Modern deep forecasting models are evaluated almost exclusively in the
enterprise-scale regime: thousands of series, years of clean daily history, and
curated preprocessing. The businesses that most need demand forecasting operate
nowhere near that regime. A small retailer typically holds six to eighteen months
of history across dozens of stock-keeping units, punctuated by stockouts,
reporting gaps and long runs of zero demand. Whether the accuracy advantages
reported for deep architectures survive this transition is, to our knowledge,
untested in a controlled way.

We introduce a \emph{scarcity sampling protocol} that turns data scarcity into a
measurable experimental factor rather than a property of whichever dataset
happens to be at hand. The protocol varies history length, series count,
missingness rate and missingness shape on an explicit grid, stratifies series by
the Syntetos--Boylan demand-pattern quadrants, and emits a manifest that makes
every cell exactly reproducible. Two design features are central: unobserved
periods are held distinct from observed zeros throughout, and series selection is
seeded independently of the scarcity axes so that comparisons across history
lengths are paired rather than confounded.

\gap{The empirical crossover results, the model comparison, and the attribution
stability analysis are not yet available. This draft presents the protocol and
its validation. No performance numbers appear anywhere in this document, and none
should be inferred.}
\end{abstract}

\section{Introduction}

Forecasting research has converged on a small number of large benchmarks. The M5
competition \citep{makridakis2022m5} provides 30{,}490 item--store series with
1{,}941 daily observations each; deep architectures such as N-BEATS
\citep{oreshkin2020nbeats} and the Temporal Fusion Transformer
\citep{lim2021tft} report their headline results on data of this scale. The
implicit claim carried by these results is that learned representations
generalize across a large pool of related series, and that the pooling is what
buys the accuracy.

That claim, if true, has an immediate corollary that the literature rarely
states: the advantage should decay as the pool shrinks. A firm with forty SKUs
and nine months of history is not a scaled-down M5. It is a qualitatively
different estimation problem, in which the number of free parameters in a
transformer may exceed the number of observations available to fit them, and in
which the classical intermittent-demand estimators of Croston
\citep{croston1972} and its successors \citep{syntetos2005accuracy,teunter2011}
were designed to operate.

Practitioners in this regime receive little guidance. The question they face is
concrete and answerable: \emph{at what data volume does a deep model stop
outperforming a classical statistical baseline, and what governs that
threshold?} Answering it requires more than re-running existing benchmarks on
subsets, because subsetting introduces confounds that are easy to miss. Shorter
windows change which series are eligible; injected missingness perturbs the very
statistics used to categorize demand patterns; and independent resampling across
grid cells discards the statistical power that paired designs provide.

\paragraph{Contributions.} This work makes the following contributions.
\begin{enumerate}
  \item A \textbf{scarcity sampling protocol} (Section~\ref{sec:protocol}) that
    treats history length, series count, missingness rate and missingness shape
    as controlled experimental factors, stratified by demand pattern, with
    deterministic manifests for exact reproduction of any grid cell.
  \item A \textbf{dual classification scheme} that separates the stable
    demand-pattern label used for stratification from the label an analyst could
    actually compute in the scarce regime, and reports the disagreement between
    them as a quantity of interest rather than as noise.
  \item An analysis of \textbf{confounds specific to scarcity experiments},
    including a missingness treatment that provably preserves the average
    inter-demand interval in expectation, and a frame-based eligibility rule that
    prevents history length from becoming entangled with series longevity.
  \item \gap{An empirical characterization of the crossover point across three
    datasets, five seeds and the full scarcity grid, together with an analysis of
    SHAP attribution stability under scarcity.}
\end{enumerate}

\section{Related Work}

\paragraph{Intermittent demand.} Croston's method \citep{croston1972} decomposes
intermittent series into demand sizes and inter-demand intervals, estimating each
separately. Syntetos and Boylan identified and corrected a bias in the original
estimator \citep{syntetos2005accuracy}, and Teunter, Syntetos and Babai proposed
the TSB variant to handle obsolescence \citep{teunter2011}. These methods remain
strong baselines precisely because they encode structural assumptions that
compensate for scarce data.

\paragraph{Demand categorization.} \citet{syntetos2005categorization} partition
series by average inter-demand interval (ADI) and the squared coefficient of
variation of demand sizes ($\mathrm{CV}^2$), yielding the smooth, erratic,
intermittent and lumpy quadrants used throughout this work. \citet{kostenko2006note}
propose an alternative boundary. We adopt the standard cutoffs, since
data-driven boundaries would make quadrants incomparable across datasets.

\paragraph{Deep forecasting.} N-BEATS \citep{oreshkin2020nbeats} and the
Temporal Fusion Transformer \citep{lim2021tft} represent two dominant families of
deep forecaster. Gradient-boosted trees \citep{ke2017lightgbm} occupy a middle
ground and were heavily represented among M5 winners
\citep{makridakis2022m5}. Our interest is not in any one architecture but in how
the ordering among these families changes with data volume.

\paragraph{Evaluation.} We follow \citet{hyndman2006another} in adopting scaled
errors, and \citet{tashman2000} and \citet{bergmeir2012} in preferring
rolling-origin evaluation to a single holdout. Pairwise comparisons use the
Diebold--Mariano test \citep{diebold1995comparing} with the small-sample
correction of \citet{harvey1997testing}, controlled across the grid by the
Benjamini--Hochberg procedure \citep{benjamini1995controlling}.

\paragraph{The gap.} Sample efficiency is studied extensively in general machine
learning but rarely in forecasting, where the dominant benchmarks fix the data
volume by construction. We are not aware of a controlled study that varies data
volume as an experimental factor while holding the series population, the
evaluation protocol and the tuning budget constant.

\section{The Scarcity Sampling Protocol}
\label{sec:protocol}

\subsection{Panel representation}

Every dataset is normalized to a long panel of triples
$(\text{series}, t, y)$ on a regular daily index. The representation makes one
distinction that the rest of the protocol depends on:
\[
y_{i,t} =
\begin{cases}
\text{NaN} & \text{the demand process was not observable at } t,\\
0 & \text{the process was observed and produced zero demand.}
\end{cases}
\]

Raw retail data conflates these. An item that launches midway through the M5
window carries a long run of leading zeros that are not demand observations at
all; the Favorita export omits rows entirely on days without a sale, so absence
is ambiguous between a closed store and genuine zero demand. Conflating the two
inflates ADI, distorts $\mathrm{CV}^2$, and pushes series into the wrong
quadrant. Because the benchmark stratifies on exactly these statistics, the error
would not stay local: it would propagate into every reported result. Loaders
therefore trim each series to its first strictly positive observation and mark
everything prior as unobserved.

\subsection{Demand-pattern stratification}

For each series we compute ADI as the number of observed periods divided by the
number of demand occurrences, and $\mathrm{CV}^2$ as the squared coefficient of
variation of the strictly positive demand sizes, and assign a quadrant using the
standard cutoffs $\mathrm{ADI} = 1.32$ and $\mathrm{CV}^2 = 0.49$
\citep{syntetos2005categorization}. Series with fewer than five demand
occurrences are marked insufficient and excluded from the sampling frame, since
$\mathrm{CV}^2$ estimated from fewer points is not an estimate of anything.

Classification is performed twice, and this is deliberate. The
\emph{frame label} is computed over full history, is fixed per series, and drives
both stratified sampling and analysis grouping. The \emph{window label} is
recomputed inside each sampled window and reflects what an analyst holding only
that window could determine. The frame label uses information outside the
sampled window, but it is never supplied to any model and therefore leaks nothing
into forecasts; it functions purely as a blocking factor, which keeps history
length and quadrant membership independent. The disagreement between the two
labels is reported per cell as the \emph{quadrant drift rate}, and is treated as
a finding in its own right.

\subsection{Missingness and the ADI invariance}
\label{sec:adi}

Missingness is injected at rates $\{0, 0.05, 0.15, 0.30\}$ under two patterns.
MCAR masks observed cells independently. Burst masks contiguous runs with
geometric lengths averaging seven days, approximating the failure modes that
actually produce gaps in small-business data: a till failing, a store closing, an
export job silently breaking for a week. Both are trimmed to the requested rate
exactly, so the patterns are compared at identical volumes of absence and differ
only in shape. Injection masks only cells that were observed, so the requested
rate carries the same meaning regardless of how much native absence a dataset
already contains.

The definition of ADI in the preceding subsection is not the conventional one,
and the reason is a confound that would otherwise be fatal to the design.
Consider counting unobserved periods as part of the inter-demand interval. Under
missingness at rate $p$, the demand count falls in expectation to $(1-p)n$ while
the period count remains $N$, giving
\[
\mathbb{E}[\widehat{\mathrm{ADI}}] \;=\; \frac{N}{(1-p)n} \;=\; \frac{\mathrm{ADI}}{1-p},
\]
an inflation of $43\%$ at $p = 0.30$. That is far more than enough to drive
series across the $1.32$ cutoff. The missingness axis of the grid would silently
become a second intermittency axis, and the two effects would be inseparable in
the results. Restricting the period count to observed periods reduces both
numerator and denominator proportionally, so
$\mathbb{E}[\widehat{\mathrm{ADI}}] = N/n = \mathrm{ADI}$ and the axes stay
orthogonal. Section~\ref{sec:validation} verifies this empirically.

\subsection{Paired seeding}

Each cell draws on two independent seeds. The \emph{selection seed} is derived
from the sampling frame alone: dataset, series count, replicate index, frame
length, coverage threshold, anchor date and allocation mode. It deliberately
excludes history length and every missingness parameter. Consequently the same
series appear in every cell along those axes, and a comparison between a 60-day
cell and a 730-day cell is a paired comparison on identical series rather than
two independent draws. This is a substantial gain in statistical power at no
cost, and it removes changes in series composition as a competing explanation for
any crossover observed. The \emph{injection seed} is derived from the complete
specification, so each cell still receives independent gaps.

\subsection{Frame-based eligibility}

A series is eligible if it has at least $80\%$ observed coverage over a fixed
898-day frame, the longest history on the grid plus the shared test period, regardless of how much history the cell subsequently hands to the
model. Judging eligibility against \texttt{history\_days} instead would let short
cells draw from a larger pool containing short-lived products while long cells
drew only from long-lived ones, entangling history length with series longevity.
The cost is some realism, since a genuinely new business does hold short-lived
series; the benefit is that an observed crossover is attributable to history
length rather than to the population shifting underneath the comparison.

\subsection{Allocation and shortfall}

Under \emph{equal} allocation each quadrant receives $n/4$ series, which
maximizes power per quadrant and is the purpose of stratifying. Under
\emph{proportional} allocation quadrants receive shares matching the eligible
pool, resolved by largest remainder. Proportional allocation exists because it is
required: as Section~\ref{sec:validation} shows, one of our datasets is almost
entirely smooth and cannot support equal allocation at all. When a quadrant holds
fewer eligible series than requested, the sampler takes what exists and records
the deficit in the manifest rather than topping up from other quadrants, which
would silently unbalance the strata the design rests on.

\begin{algorithm}[t]
\caption{Drawing one cell of the scarcity grid}
\label{alg:sample}
\begin{algorithmic}[1]
\Require dataset $D$, specification $s$, frame labels $L$
\State $\text{frame} \gets$ window of $s.\text{frame\_days}$ ending at $s.\text{end\_date}$
\State $E \gets \{i : \text{coverage}_i(\text{frame}) \ge s.\text{min\_coverage},\ L_i \ne \text{insufficient}\}$
\State $r \gets \textsc{Allocate}(s, \{|E_q|\}_{q \in Q})$ \Comment{equal or proportional}
\State $\rho_{\text{sel}} \gets \textsc{Rng}(\textsc{Hash}(\textsc{FrameKey}(s)))$ \Comment{excludes history and missingness}
\State $S \gets \bigcup_{q \in Q} \textsc{DrawWithoutReplacement}(E_q, \min(r_q, |E_q|), \rho_{\text{sel}})$
\State $W \gets$ last $s.\text{history\_days}$ of frame, restricted to $S$
\State $\rho_{\text{inj}} \gets \textsc{Rng}(\textsc{Hash}(s))$ \Comment{full specification}
\State $W \gets \textsc{InjectMissingness}(W, s.\text{rate}, s.\text{pattern}, \rho_{\text{inj}})$
\State $L^{W} \gets \textsc{Classify}(W)$ \Comment{window labels, for drift}
\State \Return $W$, manifest containing $S$, $r$, realized rates, $\textsc{Drift}(L, L^{W})$, panel digest
\end{algorithmic}
\end{algorithm}

\subsection{Reproducibility}

Each cell emits a manifest recording the full specification, both derived seeds,
window boundaries, requested and realized allocation per quadrant, any shortfall,
native and injected and final absence rates, window quadrant counts, the drift
rate, the selected series identifiers, and a SHA-256 digest of the resulting
panel. Datasets carry two content hashes, one over the raw download and one over
the normalized panel, so that an upstream file change and a change in our own
normalization code are distinguishable. Results are comparable only when both
still match, and overriding a mismatch requires an explicit flag.

\section{Experimental Setup}

\subsection{Datasets}

Three datasets span retail and non-retail demand. Table~\ref{tab:datasets}
summarizes them; only CTA has been ingested at the time of writing.

\begin{table}[t]
\centering
\caption{Datasets. Statistics for CTA are measured from the normalized panel.
M5 and Favorita figures are as documented by their sources and are not yet
independently verified in this pipeline.}
\label{tab:datasets}
\small
\begin{tabular}{llrrl}
\toprule
Dataset & Domain & Series & Span & Status \\
\midrule
M5 \citep{makridakis2022m5} & Retail, hierarchical & 30{,}490 & 1{,}941 d & Not yet ingested \\
Favorita & Grocery, perishables & $\sim$200{,}000 & 1{,}684 d & Not yet ingested \\
CTA & Rail ridership & 148 & 9{,}312 d & Ingested \\
\bottomrule
\end{tabular}
\end{table}

CTA rail station entries serve as a non-retail control. A single aggregate
ridership series would have been simpler to obtain but useless here, since the
protocol samples between 10 and 1{,}000 series; the station-level panel supports
the identical protocol and additionally contains a genuine structural break in
2020.

\subsection{Models}

Classical baselines comprise seasonal naive, ETS and auto-ARIMA
\citep{hyndman2008automatic}, with Croston \citep{croston1972} and TSB
\citep{teunter2011} for intermittent series. The machine-learning tier is
LightGBM \citep{ke2017lightgbm} over lag, rolling and calendar features. The deep
tier is N-BEATS \citep{oreshkin2020nbeats} and the Temporal Fusion Transformer
\citep{lim2021tft}. All models share a uniform fit/predict interface, an
identical feature budget and an identical tuning protocol.

\subsection{Metrics and inference}

MASE and RMSSE \citep{hyndman2006another} are primary, WAPE is secondary, and
pinball loss is reported for quantile-capable models. MAPE is not used anywhere,
as it is undefined on the zero-demand periods that dominate the intermittent
regime. Because the protocol injects gaps across the entire window including the
evaluation region, unobserved actuals are excluded from every metric rather than
coerced to zero; treating them as zeros would systematically favor whichever
model degrades most gracefully on missing inputs.

Evaluation is rolling-origin \citep{tashman2000,bergmeir2012} with at least five
replicates per configuration. Confidence intervals on metric differences are
obtained by bootstrap. Pairwise comparisons use Diebold--Mariano
\citep{diebold1995comparing} with the Harvey--Leybourne--Newbold correction
\citep{harvey1997testing}, and the false discovery rate across the grid is
controlled by Benjamini--Hochberg \citep{benjamini1995controlling}.

\subsection{Splitting: a shared evaluation period}
\label{sec:splits}

Carving rolling-origin folds out of each window separately breaks the central
comparison in two ways simultaneously.

The first is correlation. With a 28-day horizon, a 60-day window leaves roughly
32 days for training and accommodates only a couple of origins, forcing test
windows to overlap heavily; a 730-day window accommodates dozens of disjoint
folds. Overlapping test windows are correlated, which inflates the effective
sample size assumed by Diebold--Mariano tests and bootstrap intervals, and
inflates it by a different amount in every cell. The inferential machinery would
be miscalibrated inconsistently along the very axis under study.

The second is that cells would be scored on different days. A metric difference
between a 60-day cell and a 730-day cell would then confound the effect of
training volume with the effect of being graded on different data.

Both problems dissolve by moving the split outside the window. We fix a single
evaluation period in calendar time, shared by every cell: identical dates,
identical series, identical targets. History length then governs only how much
training data precedes that period. This is the standard construction of a
learning-curve experiment, and it renders the crossover literally as two curves
meeting over a shared abscissa.

The training window slides rather than expands. Under an expanding window a cell
labelled 60 days would have observed 200 days by the final fold, and the label
would be false. Sliding at fixed width holds the training size at exactly
$H$ observations in every fold. We use $h = 28$ with six folds, giving a 168-day
evaluation period; $h = 28$ matches the M5 convention, preserving
commensurability with published results, and corresponds to a monthly reorder
cycle.

One property deserves explicit statement because it superficially resembles
leakage and is not. From the second fold onward, the training window contains
days that were the preceding fold's test targets. This is precisely what
rolling-origin evaluation means \citep{tashman2000}: at each origin the
forecaster may use every actual observed before that origin, and no fold is ever
evaluated on data it was trained on.

\section{Protocol Validation}
\label{sec:validation}

No model has been fit. This section reports only measurements of the sampling
protocol itself, which is what is currently implemented and tested.

\subsection{The non-retail control has no intermittency}

Applying the classifier to the full CTA panel yields the mix in
Table~\ref{tab:cta-quadrants}: essentially every station is smooth. This is not a
defect. Rail ridership is a daily continuous process, and the classifier
correctly reports that it contains no intermittency to stratify on.

\begin{table}[h]
\centering
\caption{Syntetos--Boylan quadrants over full CTA history (148 stations,
2001-01-01 to 2026-06-30).}
\label{tab:cta-quadrants}
\small
\begin{tabular}{lrrrrr}
\toprule
& Smooth & Erratic & Intermittent & Lumpy & Insufficient \\
\midrule
Stations & 145 & 1 & 1 & 0 & 1 \\
\bottomrule
\end{tabular}
\end{table}

The consequence is methodological. Equal allocation degenerates on CTA: a request
for 20 series returns 6, because three strata are empty. This is what motivated
proportional allocation, and it repositions CTA within the benchmark. It is not a
third source of stratified evidence but a contrast case, establishing what the
crossover looks like when intermittency is absent. If deep models retain their
advantage further down the volume curve on CTA than on M5, that is direct
evidence that intermittency rather than sheer data volume is the operative
mechanism.

\subsection{Quadrant drift is an intermittency effect, not a window-length effect}

Table~\ref{tab:drift} reports drift rates. On real CTA data drift is zero at every
history length including 60 days: smooth series are classified robustly from very
little data. On synthetic panels containing all four quadrants, drift rises as
history shrinks, reaching $7.5\%$ at 60 days.

Taken together these sharpen the claim available to us. Short histories do not
mislabel series indiscriminately; they mislabel specifically those series whose
$\mathrm{CV}^2$ rests on a handful of demand occurrences, which are exactly the
series for which the choice of forecasting method matters most.

\begin{table}[h]
\centering
\caption{Quadrant drift rate, the fraction of sampled series whose quadrant
changes when estimated from the window rather than from full history. Synthetic
panels use 25 series per quadrant with cleanly separated generators and therefore
represent a lower bound; real catalogues sit closer to the cutoffs.}
\label{tab:drift}
\small
\begin{tabular}{lcccc}
\toprule
History (days) & 60 & 120 & 365 & 730 \\
\midrule
CTA, no missingness            & 0.000 & 0.000 & 0.000 & 0.000 \\
CTA, 30\% burst                & 0.000 & 0.000 & 0.000 & 0.000 \\
Synthetic, no missingness      & 0.075 & 0.025 & 0.000 & 0.000 \\
Synthetic, 30\% MCAR           & 0.075 & 0.000 & 0.000 & 0.000 \\
\bottomrule
\end{tabular}
\end{table}

\subsection{The ADI invariance holds empirically}

Section~\ref{sec:adi} argued that restricting the period count to observed
periods leaves ADI unchanged in expectation under missingness. The final two rows
of Table~\ref{tab:drift} provide the check: at 60 days of history, drift is
$0.075$ both without missingness and at $30\%$ MCAR. Injecting gaps into nearly a
third of the panel does not reshuffle quadrant membership. Under the conventional
definition the same manipulation would have inflated ADI by $43\%$.

\subsection{The pipeline reproduces the intended design on real data}

Sampling, splitting and metric evaluation were run end to end on the CTA panel at
two history lengths. Both cells score an identical 1{,}344 points over the same
evaluation period beginning 2026-01-14, while training on 480 and 5{,}840
observations respectively. Identical targets under differing training volume is
the design intent of Sections~\ref{sec:protocol} and~\ref{sec:splits}, here
confirmed against real data rather than asserted. The forecaster used was
seasonal naive, employed solely to exercise the pipeline; its accuracy is not
reported and is not a result.

\subsection{A defect visible only in real data}

The CTA feed contains 1{,}236 rows that duplicate a station--day, all within a
41-day window in mid-2011, each reporting a different ride count for the same
day. The affected share is $0.09\%$ of the panel, with median disagreement
$0.3\%$ and maximum $20\%$. No principled basis exists for preferring either
figure, and averaging them would fabricate an observation that was never
recorded, so the affected station--days are marked unobserved. We note this less
for its magnitude than for its provenance: it was invisible to a test suite
built on synthetic panels and surfaced within seconds of contact with the real
feed.

\section{Results}

\gap{This section will report the crossover analysis across the full scarcity
grid: metric surfaces by history length and series count for each model family,
stratified by demand-pattern quadrant; bootstrap confidence intervals on metric
differences; Diebold--Mariano tests with multiplicity control; and SHAP
attribution stability across seeds and scarcity regimes. None of this has been
run. Deliberately left empty rather than estimated, projected or illustrated.}

\section{Limitations}

Several limitations are known in advance and are properties of the design rather
than oversights.

\textbf{Frame-based eligibility trades realism for attribution.} Requiring 730
days of coverage excludes precisely the short-lived products a genuinely new
business would carry. The benchmark therefore measures the effect of \emph{showing
a model less history}, not the effect of \emph{a business being young}. These are
different questions and only the first is cleanly identifiable.

\textbf{Stratification uses out-of-window information.} The frame label could not
be computed by a practitioner holding only the sampled window. It is never given
to a model, so no forecast is contaminated, but the strata are in this sense
idealized. The drift rate is reported precisely to quantify what that idealization
conceals.

\textbf{The non-retail control is not stratified.} CTA is almost entirely smooth
and contributes evidence only about the smooth quadrant.

\textbf{Injected missingness is not real missingness.} MCAR and geometric bursts
are tractable idealizations. Real gaps correlate with demand itself: stockouts
suppress recorded sales precisely when demand is high, which is missingness not
at random \citep{little2019}. Our design cannot capture that dependence and will
therefore, if anything, understate the difficulty of the real problem.

\section{Conclusion}

\gap{Conclusions require results. This section is intentionally empty.}

\section*{Reproducibility}

All code, configurations and design rationale are available at
\url{https://github.com/aryanshirodkar15/scarce-forecasting}. Raw data is not
redistributed; download scripts and content hashes are provided so that any
reported figure can be traced to the exact bytes that produced it. Every grid
cell is reconstructible from its manifest.

\bibliography{refs}

\end{document}
